%==============================================================================
% 一、问题重述
%==============================================================================
\section{问题重述}

\subsection{问题背景}

啤酒的风味、口感与香气品质高度依赖其化学组成。附件数据测量了商用啤酒的化学特性，包括酒精含量、异$\alpha$-酸、pH、糖浓度等常见酿造参数，以及二百余种风味化合物。这些化合物按来源可分为酒花萜类、酵母代谢产物即酯类与醇类、麦芽来源化合物、野生菌群产物与香料或陈化产物五类；热量值、总酸、总酯、酒花香气与硫化合物五项属性由多个单独测量的化合物计算得到。

附件给出的是一套真实的风味代谢实测数据集，含有真实缺失值与异常值，且部分化合物丰度经过对数变换。因此本题的第一项工作不是拟合模型，而是把数据的量纲、缺失机制与标签尺度逐条核实清楚：只有这三件事确定之后，回归、归因与生成才有可比较的意义。

\subsection{数据概况}

附件共 8 个文件，实测结构与题面描述存在若干出入，本文一律以实测为准。核心事实为：样本 250 款啤酒、22 个风格类；化学指标 231 列，其中 226 列为逐项实测、5 列为由多项实测加和算出的派生量，含热量值、总酸、总酯、酒花香气与硫化合物，231 列全部入模——依据是附件2 提供的相关矩阵正是 $231\times231$，出题方把这 5 个派生量一并计入化合物集合；感官标签为附件4 的 50 个维度；附件3 是评分卡模板而不含真值；附件8 提供化学族分组、香气阈值与对数变换对照表。详细核对见第五章。

\subsection{问题提出}

问题一：数据缺失条件下的感官评分。以附件1 化合物数据为基础，给出附件3 中定义的 Aroma 与 Flavour 两套、共约 40 个感官维度的 0--6 分制分级评分结果。

问题二：化合物--风味分析。利用附件2 的化合物相关数据构建化合物共表达图，结合附件1 特征，给出每个感官维度的 Top-3 关键因素化合物，并进行稳定性与化学合理性论证。

问题三：风味与“风味指纹”生成。以 tasting\_category\_fine 为风格标签学习风格条件生成模型，在给定风格下生成合理的化合物谱并给出其感官画像；同时做跨风格迁移，并保持 DMS、Diacetyl、4VG、trans-2-nonenal 等缺陷化合物不恶化。
